\clearpage
\SourceLeaf{34}

\begin{center}\large A NEW FORM OF THE PROOF OF THE IRRATIONALITY OF THE TANGENTS OF RATIONAL ARCS.\end{center}

\noindent\hspace*{\parindent}The infinite series

\[
P = 1-\frac{1}{2} . \frac{mm}{nn}+\frac{1}{2.4} . \frac{1}{1.3} . \frac{m^4}{n^4}-\frac{1}{2.4.6} . \frac{1}{1.3.5} . \frac{m^6}{n^6}+\text{and\,so\,on}
\]
\[
\begin{aligned}
P_1 &= \frac{m}{n}-\frac{1}{2} . \frac{1}{1.3} . \frac{m^2}{n^3}\\ &\quad+\frac{1}{2.4} . \frac{1}{1.3.5} . \frac{m^5}{n^5}-\frac{1}{2.4.6} . \frac{1}{1.3.5.7} . \frac{m^7}{n^7}+\text{and\,so\,on}
\end{aligned}
\]
\[
\begin{aligned}
P_2 &= \frac{1}{1.3} . \frac{m^3}{n.n}-\frac{1}{2} . \frac{1}{1.3.5} . \frac{m^5}{n^4}\\ &\quad+\frac{1}{2.4} . \frac{1}{1.3.5.7} . \frac{m^7}{n^6}-\frac{1}{2.4.6} . \frac{1}{1.3.5.7.9} . \frac{m^9}{n^8}+\text{and\,so\,on}
\end{aligned}
\]
\[
\begin{aligned}
P_3 &= \frac{1}{1.3.5} . \frac{m^5}{n^3}-\frac{1}{2} . \frac{1}{1.3.5.7} . \frac{m^7}{n^5}\\ &\quad+\frac{1}{2.4} . \frac{1}{1.3.5.7.9} . \frac{m^9}{n^7}-\frac{1}{2.4.6} . \frac{1}{1.3.5.7.9.11} . \frac{m^{11}}{n^9}+\text{and\,so\,on}
\end{aligned}
\]

\noindent and in general, starting from $\theta = 1$,

\[
\begin{aligned}
P_\theta = {}& \frac{1}{1.3.5\ldots(2\theta-1)} . \frac{m^{2\theta-1}}{n^\theta}
-\frac{1}{2} . \frac{1}{1.3.5\ldots(2\theta+1)} . \frac{m^{2\theta+1}}{n^{\theta+2}}
\\ &\,+\frac{1}{2.4} . \frac{1}{1.3.5\ldots(2\theta+3)} . \frac{m^{2\theta+3}}{n^{\theta+4}}
\\ &\,-\frac{1}{2.4.6} . \frac{1}{1.3.5\ldots(2\theta+5)} . \frac{m^{2\theta+5}}{n^{\theta+6}}+\text{and\,so\,on},
\end{aligned}
\]

\noindent where $m$, $n$ initially denote arbitrary given positive quantities,
all evidently converge, if only after being carried sufficiently far,
and converge faster than any decreasing geometric progression.
Each therefore has a definite finite numerical value.

If the second term of $P_\theta$ is already smaller than the first, or even
equal to it, the third is automatically smaller than the second, the
fourth smaller than the third, and\,so\,on. This case requires only
that $2\theta+1$ be at least $\dfrac{mm}{2nn}$. Then the
value of $P_\theta$ is certainly positive, smaller than the first term, and
greater than the sum of the first two.
